\documentclass[11pt]{article}

% ---------------------------------------------------------------------------
% Portable preamble: compiles with pdfLaTeX or XeLaTeX, no project macros.
% Preamble macros copied verbatim from theory.tex (the companion fixed-grid
% document) for notational consistency, plus three new macros for the
% row-$n$ discretized nuisances used throughout this document.
% ---------------------------------------------------------------------------
\usepackage{amsmath,amsthm,amssymb}
\usepackage[margin=1in]{geometry}
\usepackage[colorlinks=true,linkcolor=blue,urlcolor=blue,citecolor=blue]{hyperref}
\usepackage{mathtools}
\usepackage{enumitem}
\usepackage{booktabs}
\usepackage{longtable}
\usepackage[round,authoryear]{natbib}

\theoremstyle{plain}
\newtheorem{theorem}{Theorem}[section]
\newtheorem{lemma}[theorem]{Lemma}
\newtheorem{proposition}[theorem]{Proposition}
\newtheorem{corollary}[theorem]{Corollary}
\theoremstyle{definition}
\newtheorem{assumption}{Assumption}
\newtheorem{definition}[theorem]{Definition}
\theoremstyle{remark}
\newtheorem{remark}[theorem]{Remark}
\newtheorem{example}[theorem]{Example}

\renewcommand{\theassumption}{B\arabic{assumption}}

\newcommand{\E}{\mathbb{E}}
\newcommand{\PP}{\mathbb{P}}
\newcommand{\Pn}{\mathbb{P}_n}
\newcommand{\R}{\mathbb{R}}
\newcommand{\cX}{\mathcal{X}}
\newcommand{\cA}{\mathcal{A}}
\newcommand{\cT}{\mathcal{T}}
\newcommand{\cS}{\mathcal{S}}
\newcommand{\cP}{\mathcal{P}}
\newcommand{\cL}{\mathcal{L}}
\newcommand{\cG}{\mathcal{G}}
\newcommand{\ez}{e_0}
\newcommand{\mz}{\mu_0}
\newcommand{\wz}{w_0}
\newcommand{\ehat}{\widehat{e}}
\newcommand{\mhat}{\widehat{\mu}}
\newcommand{\what}{\widehat{w}}
\newcommand{\that}{\widehat{\tau}}
\newcommand{\thetahat}{\widehat{\theta}}
\newcommand{\thetatil}{\widetilde{\theta}}
\newcommand{\pihat}{\widehat{\pi}}
\newcommand{\Vhat}{\widehat{V}}
\newcommand{\Dw}{D_w}
\newcommand{\Dm}{D_\mu}
\newcommand{\Ltwo}[1]{\lVert #1 \rVert_{2}}
\newcommand{\Lw}[1]{\lVert #1 \rVert_{2,w}}
\newcommand{\Linf}[1]{\lVert #1 \rVert_{\infty}}
\DeclareMathOperator*{\argmin}{arg\,min}
\DeclareMathOperator*{\argmax}{arg\,max}
\DeclareMathOperator{\Var}{Var}
\DeclareMathOperator{\plim}{plim}
\DeclarePairedDelimiter{\abs}{\lvert}{\rvert}
\DeclarePairedDelimiter{\card}{\lvert}{\rvert}

%% Row-$n$ discretized nuisances: \ez, \mz, \wz already expand to e_0, mu_0,
%% w_0 (the FIXED-GRID truth, as in the companion document); appending a
%% second subscript directly (\ez_n) produces a double-subscript TeX error.
%% \ezn, \mzn denote the row-$n$ discretized nuisances e_{0,n}, mu_{0,n} of
%% \eqref{eq:mesh-en}/\eqref{eq:mesh-mun} -- distinct objects from \ez/\mz,
%% not re-subscripted notation for the same symbol. The CONTINUUM truth is
%% written \ez^\circ, \mz^\circ throughout (no new macro needed).
\newcommand{\ezn}{e_{0,n}}
\newcommand{\mzn}{\mu_{0,n}}

\title{Inference for Optimal-Tree Nuisances Under a Refining
Discretization Grid\\[2pt]
\large A companion to the fixed-grid theory}
\author{}
\date{}

\begin{document}
\maketitle

\begin{abstract}
We study inference for the average treatment effect on the treated when both
nuisance functions are fitted by a globally optimal decision tree with a
leaf budget $\bar L$ that is fixed, but whose covariate discretization grid
is allowed to refine as the sample size grows: the analyst's grid $\cA_n$ is
generated by a discretization map $X_n=Q_n(Z)$ of an underlying continuum
covariate $Z$, with mesh width $h_n\to0$. This is a distinct regime from
either growing the leaf budget or dropping the pre-specified-cutpoint
restriction entirely, both of which are excluded here and developed
elsewhere. We give two results. First, a population-level approximation
lemma: the grid-alignment error between a continuum nuisance and its best
grid-restricted tree approximation vanishes at rate $O(\sqrt{h_n})$, so that
exact leaf-constancy of the true nuisances is required only at the
continuum level, not at the analyst's particular choice of grid. Second, an
inferential half built on a genuine oracle inequality --- replacing exact
selection consistency, which provably cannot survive grid refinement under
any uniformity assumption --- that delivers: a cross-fitted estimator that is
$\sqrt n$-consistent, asymptotically normal, and efficient whenever
$h_n=o(n^{-1/2})$; a conservative anchor interval whose width contracts at a
stated, possibly slower, rate; and, under an additional joint
empirical-process condition, a full-sample estimator that is itself
$\sqrt n$-consistent, asymptotically normal, and locally efficient. This
document depends on, but does not modify, a companion fixed-grid theory
document, from which every reused assumption and background result is
restated here in full; proofs of those background results are not
reproduced and are cited to the companion document.
\end{abstract}

\tableofcontents

% ===========================================================================
\section{Introduction and scope}\label{sec:intro}

The companion fixed-grid theory develops $\sqrt n$-efficient inference for
the average treatment effect on the treated when both nuisance functions are
fitted by a single globally optimal decision tree, with a fixed leaf budget
$\bar L$ over a fixed, finite candidate class of grid atoms $\cA$. That
theory's efficiency result requires \emph{structural sparsity}: that the true
propensity and outcome regression are each exactly constant on the leaves of
some tree in the candidate class. A natural objection is that this asks the
truth to align with the analyst's particular choice of grid, which is a
property of the discretization, not of the underlying phenomenon.

This document gives the formal answer, in two parts, for the regime in which
the leaf budget $\bar L$ stays exactly fixed while the grid itself refines:
$\cA_n$ is generated by a discretization map $X_n=Q_n(Z)$ of an underlying
continuum covariate $Z$, with mesh width $h_n\to0$.

\emph{Part I} (Section~\ref{sec:mesh}) is population-level and
approximation-theoretic: it shows that the grid-alignment bias between a
continuum nuisance and the best tree the analyst's grid can represent
vanishes at rate $O(\sqrt{h_n})$, so exact sufficiency is required only of
the continuum-level nuisances, never of the analyst's particular
discretization.

\emph{Part II} (Section~\ref{sec:grid-inference}) is inferential. Grid
refinement breaks the fixed-grid theory's route to exact selection
consistency in a way that cannot be repaired by any uniformity assumption
(Section~\ref{sec:grid-inference} states exactly why), so a genuine
finite-class oracle inequality is developed to replace it. Building on that
oracle inequality, three results are given, with deliberately different and
explicitly stated strength: a cross-fitted estimator that is a genuine
corollary of the companion document's cross-fitted central limit theorem; an
anchor-interval width result that is honestly a \emph{rate}, not an
\emph{efficiency}, statement; and a full-sample estimator's central limit
theorem that is an independent result, not a corollary of anything in the
companion document, since it must be established without the sparsity
condition the companion document's main efficiency theorem relies on.

\paragraph{Scope exclusions.} Three related regimes are explicitly
\emph{not} addressed here. Growing the leaf budget, $\bar L_n\to\infty$,
reopens a uniform-in-tree-class entropy argument this document is designed
to avoid, and is developed elsewhere. Growing the raw covariate dimension,
$d_n\to\infty$, is excluded throughout: the overlap constant and the
existence of a stable continuum limit for $\cX_n$ are not guaranteed to
survive that direction. Dropping the pre-specified-cutpoint restriction
altogether (so that the tree search is not confined to grid-restricted
partitions) is a third, distinct generalization not addressed here.

\paragraph{Relationship to the companion document.} This document is
self-contained in its assumptions: every assumption it uses is restated in
full in Section~\ref{sec:background}, not invoked by reference to an
external file. It is \emph{not} self-contained in its proofs of background
material: Section~\ref{sec:background} states, but does not reprove, the
handful of results from the companion fixed-grid theory that the new
material here builds on. Every such result is clearly marked and cited to
the companion document for its proof. Everything new --- the population
approximation lemma of Section~\ref{sec:mesh} and the full inferential
apparatus of Section~\ref{sec:grid-inference} --- is proved here in full.

% ===========================================================================
\section{Setup and notation}\label{sec:setup}

Data are $n$ i.i.d.\ copies of $O=(X,A,Y)$, where $X$ is the analyst's
discretized covariate, $A\in\{0,1\}$ a binary treatment, and $Y$ an observed
outcome. Write $\ez(x)=\PP(A=1\mid X=x)$, $\mz(x)=\E(Y\mid A=0,X=x)$,
$\pi=\PP(A=1)$, $m_1(x)=\E(Y\mid A=1,X=x)$, and $\wz=\ez/(1-\ez)$. The
target parameter is the average treatment effect on the treated,
\[
  \theta_0=\pi^{-1}\,\E\bigl[\ez\,(m_1-\mz)\bigr],
\]
identified under consistency and ignorability as in
Assumption~\ref{ass:causal} below. Under the refining-grid regime of
Section~\ref{sec:grid-inference}, this display is superseded: $\theta_0$ is
restated there at the level of the continuum covariate $Z$, and
Assumption~\ref{ass:mesh} resolves explicitly why that restatement is
needed and why it is compatible with, rather than a contradiction of,
identification given $X$. The score is
\[
  \psi(O;\theta,\eta)
  =\frac1\pi\Bigl[A\{Y-\mu(X)-\theta\}-(1-A)w(X)\{Y-\mu(X)\}\Bigr],
  \qquad \eta=(e,\mu),\ w=e/(1-e),
\]
with $\psi_0=\psi(\cdot;\theta_0,\eta_0)$.

A tree partition $\tau$ of the covariate space with leaves
$\ell_1,\dots,\ell_K$ induces, for a weight $\nu$ mutually absolutely
continuous with $P$, the leaf-refit projection
$\Pi^\nu_\tau\gamma$, constant on each leaf at the $\nu$-weighted within-leaf
average of $\gamma$; write $\Pi_\tau$ for the unweighted case ($\nu=P$). The
sample analogue is
\begin{equation}\label{eq:refit}
  \mhat_\tau\restriction_\ell
  =\frac{\sum_{i:X_i\in\ell,\,A_i=0}Y_i}
        {\card{\{i:X_i\in\ell,\,A_i=0\}}},
\end{equation}
the within-leaf treated fraction for the propensity and the within-leaf
control mean for the outcome. Structure is selected by penalized empirical
risk minimization over the \emph{feasible} candidates,
\begin{equation}\label{eq:feasible}
  \cT^{(j)}_n=\bigl\{\tau\in\cT_{\bar L}:\text{every leaf of }\tau\text{
  carries at least }m_n\text{ observations, control observations for
  }j=\mu\bigr\},
\end{equation}
and, for a penalty $\lambda_n>0$,
\begin{equation}\label{eq:select}
  \that_j\in\argmin_{\tau\in\cT^{(j)}_n}
  \bigl\{R^{(j)}_n(\tau)+\lambda_n\card\tau\bigr\},
  \qquad j\in\{e,\mu\},
\end{equation}
where $R^{(j)}_n$ is the empirical risk of the leaf-refit fit for nuisance
$j$ and $R^{(j)}$ its population counterpart. We write
$\ehat=\ehat_{\that_e}$, $\mhat=\mhat_{\that_\mu}$, $\what=\ehat/(1-\ehat)$.

For the refining-grid regime of this document, an underlying continuum
covariate $Z$ is discretized by a row-indexed map $X_n=Q_n(Z)$ into a grid
$\cA_n$; the exact construction is given in Assumption~\ref{ass:mesh}. The
continuum truth is written $\ez^\circ(Z)=\PP(A=1\mid Z)$,
$\mz^\circ(Z)=\E(Y\mid A=0,Z)$, and the row-$n$ discretized truth
$\ezn,\mzn$ (its population projection onto $\sigma(X_n)$) is defined in
\eqref{eq:mesh-en}--\eqref{eq:mesh-mun}. For $j=e$ write $\|f\|_e^2=\E\{f(Z)^2\}$
and for $j=\mu$ write $\|f\|_\mu^2=\E[\{1-\ez^\circ(Z)\}f(Z)^2]$; for a
row-$n$ fitted nuisance,
\[
  D_{e,n}=\|\widehat e_n-\ez^\circ\|_e,
  \qquad
  D_{\mu,n}=\|\widehat\mu_n-\mz^\circ\|_\mu,
  \qquad
  D_{w,n}=\|\widehat w_n-w_0^\circ\|_\mu,
\]
where $w_0^\circ=\ez^\circ/(1-\ez^\circ)$.

% ===========================================================================
\section{Background: results from the fixed-grid theory}\label{sec:background}

This section restates, in full, every assumption and background result the
new material of Sections~\ref{sec:mesh}--\ref{sec:grid-inference} depends
on. Nothing in this section is proved here; the substance is developed in
the companion fixed-grid theory document (henceforth \emph{the companion
document}), to which every proof below is cited. Two statements
(Corollary~\ref{cor:width}, Theorem~\ref{thm:regular}) are restated with
added commentary or in a form marked in situ as a licensed generalization
of the companion's own version, rather than as a verbatim transcription.
The purpose of restating these here, rather than citing them by reference
alone, is that this document is meant to be read on its own.

\subsection{Standing assumptions (reused, unmodified)}\label{sec:background-assumptions}

\begin{assumption}[Solver fidelity]\label{ass:solver}
The optimisation procedure computing~\eqref{eq:select} considers only tree
partitions whose splits lie at pre-specified cutpoints, never a
data-dependent or continuously-refined set of thresholds: the search space
is $\cT_{\bar L}$ (row-indexed as $\cT_{\bar L,n}$ under
Assumption~\ref{ass:mesh}), by construction. This is a standing assumption,
in force throughout and never re-cited as a separate hypothesis below
(the companion document re-cites it explicitly in some theorem statements;
here it is suppressed throughout, since every search space appearing in any
result here is by definition built from pre-specified cutpoints).
\end{assumption}

\begin{assumption}[Sampling]\label{ass:iid}
$O_1,\dots,O_n$ are independent and identically distributed copies of
$O=(X,A,Y)\sim P$.
\end{assumption}

\begin{assumption}[Identification and one-sided positivity]\label{ass:causal}
Consistency, $Y=AY(1)+(1-A)Y(0)$; ignorability,
$\{Y(0),Y(1)\}\perp A\mid X$; $\pi>0$; and there exists $c\in(0,1)$ with
$1-\ez(x)\ge c$ for every $x\in\cX$.
\end{assumption}

\begin{assumption}[Moments and non-degeneracy]\label{ass:moments}
$\E[Y^2\mid A=a,X=x]\le C_Y<\infty$ for all $a\in\{0,1\}$ and $x\in\cX$, and
$V>0$, the semiparametric efficiency bound given explicitly in
\eqref{eq:vdecomp} below.
\end{assumption}

\begin{assumption}[Estimator construction]\label{ass:construct}
The nuisances are fitted as follows.
\begin{enumerate}[label=(\alph*),nosep]
  \item \emph{Loss.} Both nuisances are fitted by squared error, so that
  excess risk equals squared $L_2$ distance to the truth.
  \item \emph{Leaf refit.} Leaf values are given by~\eqref{eq:refit}.
  \item \emph{Clipping.} Fitted propensities are clipped so that
  $1-\ehat(x)\ge c$ for every $x$, with $c$ the constant of
  Assumption~\ref{ass:causal}.
  \item \emph{Leaf mass.} Selection is restricted to the feasible set
  \eqref{eq:feasible}: candidates whose every leaf carries at least $m_n$
  observations, control observations in the case of $\mu$, with
  $\bar Lm_n\lesssim n$.
\end{enumerate}
\end{assumption}

\begin{assumption}[Exact global optimisation]\label{ass:global}
For each $j$, $\that_j$ attains the global minimum in~\eqref{eq:select}
over the feasible set $\cT^{(j)}_n$ of~\eqref{eq:feasible}, and
$\lambda_n\to0$.
\end{assumption}

\begin{assumption}[Rate condition]\label{ass:rate}
$\Ltwo{\ehat-\ez}\cdot\Lw{\mhat-\mz}=o_p(n^{-1/2})$, where the norms are
evaluated at the fold-wise fits (fixed $K$-fold cross-fitting); this is the
product-rate condition of the debiased-machine-learning framework
\citep{chernozhukov2018}.
\end{assumption}

The following two assumptions are restated \emph{only} because
Corollary~\ref{cor:width} below cites them; they are exactly the two
fixed-grid finiteness clauses that Assumption~\ref{ass:mesh} (in
Section~\ref{sec:mesh}) replaces, and no other result in this document
depends on them.

\begin{assumption}[Finite candidate partition family]\label{ass:finite}
The grid atoms $\cA$ generated by the pre-specified cutpoints are finite in
number, fixed in $n$, and each of positive mass.
\end{assumption}

\begin{assumption}[Fixed leaf budget]\label{ass:budget}
$\bar L$ is a constant, fixed in $n$, with $\bar L\le\card\cA$;
consequently $\cT_{\bar L}$ is a fixed finite set and $\card{\cT_{\bar L}}$
does not depend on $n$.
\end{assumption}

\subsection{Background results (statements only; proofs in the companion document)}\label{sec:background-results}

\begin{proposition}[Exact bilinear remainder; companion document,
Proposition~\emph{prop:bilinear}]\label{prop:bilinear}
Under Assumptions~\ref{ass:iid}--\ref{ass:causal}, for any $\eta=(e,\mu)$
with $1-e\ge c$,
\begin{equation}\label{eq:bilinear}
  \E\bigl[\psi(O;\theta_0,\eta)\bigr]
  =\frac1\pi\,\E\Bigl[\bigl\{1-\ez(X)\bigr\}
   \bigl\{\wz(X)-w(X)\bigr\}\bigl\{\mz(X)-\mu(X)\bigr\}\Bigr].
\end{equation}
\end{proposition}

\begin{lemma}[Bias bound in the control-weighted norm; companion document,
Lemma~\emph{lem:biasbound}]\label{lem:biasbound}
Under Assumptions~\ref{ass:iid}--\ref{ass:causal} and
Assumption~\ref{ass:construct}(c), evaluated at $\eta=\widehat\eta$,
\[
  \bigl\lvert\E[\psi(O;\theta_0,\eta)]\bigr\rvert_{\eta=\widehat\eta}
  \ \le\ \frac1\pi\,\Dw\,\Dm .
\]
\end{lemma}

\begin{corollary}[Unweighted forms; companion document,
Corollary~\emph{cor:convert}]\label{cor:convert}
Under the conditions of Lemma~\ref{lem:biasbound},
\begin{equation}\label{eq:convert}
  \bigl\lvert\E[\psi(O;\theta_0,\eta)]\bigr\rvert_{\eta=\widehat\eta}
  \ \le\ \frac1{\pi c^{3/2}}\,\Ltwo{\ehat-\ez}\,\Lw{\mhat-\mz}.
\end{equation}
\end{corollary}

\begin{lemma}[Norm equivalence; companion document, Lemma~\emph{lem:normequiv}]
\label{lem:normequiv}
Under Assumption~\ref{ass:causal}, for every $f$ with $\Ltwo f<\infty$,
\[
  \sqrt c\,\Ltwo f\ \le\ \Lw f\ \le\ \Ltwo f .
\]
\end{lemma}

\begin{lemma}[One-sided clipping is non-expansive; companion document,
Lemma~\emph{lem:clip}]\label{lem:clip}
Under Assumption~\ref{ass:causal}, clipping to $[0,1-c]$ can only decrease
the error of a fitted propensity, pointwise and hence in every $L_p$ norm.
\end{lemma}

\begin{lemma}[Skeleton identity; companion document, Lemma~\emph{lem:skeleton}]
\label{lem:skeleton}
On the event $\{\pihat>0\}$, and with no expansion or approximation,
\begin{equation}\label{eq:skeleton}
  \pihat\,(\thetahat-\theta_0)=\Pn\phi(\cdot;\eta_0)+U-T,
\end{equation}
where $\phi(O;\eta)=A\{Y-\mu(X)-\theta_0\}-(1-A)w(X)\{Y-\mu(X)\}$,
$\phi(\cdot;\eta_0)=\pi\psi_0$,
\[
  U:=\Pn\bigl[h(O)\,\Delta\mu(X)\bigr],\quad
  h(O):=\wz(X)(1-A)-A,
  \qquad
  T:=\Pn\bigl[(1-A)\,\Delta w(X)\{Y-\mhat(X)\}\bigr],
\]
$\Delta\mu=\mhat-\mz$, $\Delta w=\what-\wz$, and $\E[h\mid X]=0$ pointwise.
\end{lemma}

\begin{theorem}[Cross-fitted central limit theorem; companion document,
Theorem~\emph{thm:crossfit}]\label{thm:crossfit}
Under Assumptions~\ref{ass:iid}--\ref{ass:moments},
\ref{ass:construct}, \ref{ass:global}, and~\ref{ass:rate}, with $K$-fold
cross-fitting and $K$ fixed,
\[
  \sqrt n\bigl(\thetahat_{\mathrm{cf}}-\theta_0\bigr)\rightsquigarrow N(0,V).
\]
Its proof requires no entropy or complexity condition on the candidate
class $\cT_{\bar L}$: splitting removes the need for uniform control over
the class entirely.
\end{theorem}

\begin{theorem}[Anchor interval; companion document, Theorem~\emph{thm:anchor}]
\label{thm:anchor}
Let $\thetatil$ satisfy $\sqrt n(\thetatil-\theta_0)\rightsquigarrow
N(0,\widetilde V)$ with $\widetilde V>0$ and $\widetilde\sigma^2\to_p
\widetilde V$, and let $\widehat\delta=\thetahat-\thetatil$. Then
\[
  C_n=\Bigl[\thetahat\pm\bigl(\abs{\widehat\delta}
   +z_{1-\alpha/2}\,\widetilde\sigma/\sqrt n\bigr)\Bigr]
\]
satisfies $\liminf_n\PP(\theta_0\in C_n)\ge1-\alpha$ on $\{\pihat>0\}$, with
no condition on the approximation errors and no smoothness condition on the
nuisances. Its proof is a bare triangle inequality: no property of
$\widehat\delta$ beyond its being a real number is used, so it never
produces a limit law for $\thetahat$ on its own. (A bias-aware critical
value \citep{armstrong2018}, exploiting that a bias of known magnitude
costs less than its full value in half-width, improves on this crude
construction in practice; not used below.)
\end{theorem}

\begin{corollary}[Width; companion document, Corollary~\emph{cor:width}]
\label{cor:width}
Suppose the conditions of Theorem~\ref{thm:anchor} hold, together with
Assumptions~\ref{ass:moments}--\ref{ass:construct}, \ref{ass:finite},
\ref{ass:budget}, and the conclusion of Theorem~\ref{thm:crossfit} for the
anchor. Then
\[
  \abs{\widehat\delta}=O_p\bigl(\Dw\Dm+n^{-1/2}\bigr),
  \qquad
  \mathrm{width}(C_n)=O_p\bigl(\max\{n^{-1/2},\ \Dw\Dm\}\bigr).
\]
This bound is non-asymptotic in the approximation errors $\Dw,\Dm$: no step
of its proof requires them bounded away from zero or invariant in $n$, and
the plug-in remainder in its proof is stated as a union bound over the
finite class $\cT_{\bar L}\times\cT_{\bar L}$ (Section~\ref{sec:grid-inference}
below is precisely about what happens to that union bound when the class
grows with $n$). Of the two just-added hypotheses, \ref{ass:finite}'s that
$\cA$ is fixed in $n$ and \ref{ass:budget}'s consequence that
$\card{\cT_{\bar L}}$ is $n$-invariant are exactly the two clauses
Assumption~\ref{ass:mesh} replaces; Corollary~\ref{cor:grid-width} below
therefore re-derives, rather than invokes, this result.
\end{corollary}

\begin{lemma}[Local saturation and the variance bound; companion document,
Lemma~\emph{lem:tangent}]\label{lem:tangent}
Let $\cP_{c,C_Y}$ be the set of observed-data laws of $(X,A,Y)$ on an
arbitrary measurable covariate space, satisfying $\pi>0$, $1-\ez\ge c$
$P$-a.s., $\E[Y^2\mid A=a,X=x]\le C_Y$, and $V>0$. Let $P_0$ be interior to
$\cP_{c,C_Y}$ in the sense $\ez>0$ $P_0$-a.s.\ and
$\operatorname*{ess\,inf}_{P_0}(1-\ez)>c$. Then the tangent space of
$\cP_{c,C_Y}$ at $P_0$ is the entire nonparametric tangent space, $\theta$
is pathwise differentiable at $P_0$ with canonical gradient $\psi_0$, and
the semiparametric efficiency bound is $V=\E[\psi_0^2]$, given explicitly by
\begin{equation}\label{eq:vdecomp}
  V=\pi^{-2}\,\E\Bigl[\ez\sigma_1^2
   +\frac{\ez^2\sigma_0^2}{1-\ez}
   +\ez\,(m_1-\mz-\theta_0)^2\Bigr].
\end{equation}
\end{lemma}

\begin{theorem}[Regularity and local asymptotic efficiency; companion
document, Theorem~\emph{thm:regular}]\label{thm:regular}
Let $P_0\in\cP_{c,C_Y}$ satisfy the interiority condition of
Lemma~\ref{lem:tangent}, and let $\widetilde\theta$ be \emph{any} estimator
satisfying $\sqrt n(\widetilde\theta-\theta_0)=\sqrt nP_n\psi_0+o_{P_0}(1)$
at $P_0$, for $\psi_0$ the canonical gradient of Lemma~\ref{lem:tangent}.
Then $\widetilde\theta$ is regular at $P_0$ relative to $\cP_{c,C_Y}$ and,
by the convolution theorem, locally asymptotically efficient. Its proof
(contiguity via Assumptions~\ref{ass:iid} and~\ref{ass:moments}, Le Cam's
first and third lemmas, and Lemma~\ref{lem:tangent}'s pathwise
differentiability) consumes no property of the estimator beyond this one
asymptotic-linearity statement, established at $P_0$ only: it does not use
structural sparsity, a specific estimator construction, or the
candidate-class cardinality. This is a mild generalization, in statement
only, of the companion document's own proof of this theorem for the
displayed-tree estimator $\widehat\theta$ under Theorem~\emph{thm:main}'s
hypotheses --- $\widehat\theta$'s asymptotic linearity there is simply one
instance of the hypothesis stated here.
\end{theorem}

\begin{remark}[Fixed-misspecification variance failure; companion document,
Appendix~\emph{app:honest}]\label{rem:honest}
Away from exact sufficiency, with $a=\mu_*-\mz$ and $b=w_*-\wz$ denoting
\emph{fixed}, $n$-invariant discrepancies between the leaf-projected and
true nuisances, the plug-in variance estimator satisfies
\[
  \plim\Vhat
  =V+2\pi^{-2}\bigl\{\E[\{1-\ez\}\wz b\,\sigma_0^2]
   -\E[\ez\,a\,\Delta m]\bigr\}+O\bigl(\Ltwo a^2+\Ltwo b^2\bigr),
\]
whose first-order term is not sign-definite: it is proved in the companion
document that $\plim\Vhat<V$ for a range of misspecification directions.
This is stated here only for contrast: Section~\ref{sec:grid-inference}'s
Theorem~\ref{thm:grid-fullsample} recovers $\plim\Vhat=V$ precisely because
its discrepancies $a,b$ \emph{vanish} with $n$ rather than being fixed, so
the remainder above vanishes along with the first-order term rather than
being dominated by it.
\end{remark}

% ===========================================================================
\section{Population-level approximation under grid refinement}\label{sec:mesh}

The leaf budget $\bar L$ remains fixed throughout this section; only the
covariate discretization refines. The question answered here is whether the
sufficient class depends on where the analyst happened to put the grid, not
whether the fixed-budget approximation bias can always be made to vanish.
The latter statement is false and is not claimed.

\begin{assumption}[Continuum-anchored mesh refinement]\label{ass:mesh}
Let $Z$ take values in a compact rectangle
$\cX^\circ=\prod_{r=1}^{d}[a_r,b_r]\subset\R^d$, where the dimension $d$ is
fixed. For each $n$, let $Q_n$ be a rectangular discretisation map with
\[
  X_n=Q_n(Z)\in\cX_n=\prod_{r=1}^{d}\cX_{n,r},
  \qquad
  J_{n,\max}=\max_{r\le d}\card{\cX_{n,r}},
\]
such that the largest coordinate-wise diameter of any discretisation cell is
at most $h_n$, where $h_n\longrightarrow0$. The grids are nested and
generate the Borel $\sigma$-field of $\cX^\circ$. Write
\[
  \cT_{\bar L,n}
  :=\{\text{grid-restricted tree partitions of $\cX_n$ with at most
  $\bar L$ leaves}\}
\]
for the row-$n$ candidate class, with $\bar L$ fixed throughout.

There is a constant $C_X<\infty$ such that, for every coordinate $r$, every
$t\in[a_r,b_r]$, and all sufficiently small $u>0$,
\begin{equation}\label{eq:mesh-slab}
  \PP\{\abs{Z_r-t}\le u\}\le C_Xu .
\end{equation}
The continuum nuisance functions
\[
  \ez^\circ(z)=\PP(A=1\mid Z=z),
  \qquad
  \mz^\circ(z)=\E[Y\mid A=0,Z=z]
\]
are measurable and bounded, $\abs{Y}\le B_Y<\infty$ almost surely, and
$1-\ez^\circ(z)\ge c$ for $P_Z$-almost every $z$, with the same $c$ as in
Assumption~\ref{ass:causal}. Their discretised versions are
\begin{align}
  \ezn(x)
  &=\E\{\ez^\circ(Z)\mid X_n=x\},
  \label{eq:mesh-en}\\
  \mzn(x)
  &=\frac{\E[\{1-\ez^\circ(Z)\}\mz^\circ(Z)\mid X_n=x]}
          {\E[1-\ez^\circ(Z)\mid X_n=x]} .
  \label{eq:mesh-mun}
\end{align}
These are the row-$n$ discretised versions of the continuum nuisances.

The causal model and target of every result in Section~\ref{sec:grid-inference}
are stated at the level of $Z$, not $X_n$, and this paragraph is the
resolution promised in Section~\ref{sec:setup}: consistency,
$Y=AY(1)+(1-A)Y(0)$, and ignorability, $\{Y(0),Y(1)\}\perp A\mid Z$, are
assumed (the $Z$-indexed analogue of Assumption~\ref{ass:causal});
$\E[Y^2\mid A=a,Z=z]\le C_Y$ for both $a\in\{0,1\}$ and $P_Z$-a.e.\ $z$
(implied here by $|Y|\le B_Y$ with $C_Y=B_Y^2$); and $V>0$, where, for the
remainder of this document, $V$ denotes the continuum analogue
of~\eqref{eq:vdecomp}, obtained by replacing $\ez,\mz,m_1$ throughout by
$\ez^\circ,\mz^\circ$, and $m_1^\circ(z):=\E(Y\mid A=1,Z=z)$ --- \emph{not}
the $X$-indexed quantity Assumption~\ref{ass:moments} in
Section~\ref{sec:background-assumptions} refers to; every "$V$" in
Section~\ref{sec:grid-inference} below is this continuum quantity. The
target of every result in Section~\ref{sec:grid-inference} is accordingly
\[
  \theta_0:=\pi^{-1}\,\E\bigl[\ez^\circ(m_1^\circ-\mz^\circ)\bigr],
\]
identified from consistency and ignorability given $Z$ exactly as in the
discussion following Assumption~\ref{ass:causal}, \emph{not} the
discretised, $X_n$-indexed quantity that a literal reading of
Assumptions~\ref{ass:causal}--\ref{ass:moments} in
Section~\ref{sec:background-assumptions} would produce. Ignorability given
$Z$ does \emph{not} imply, and is not implied by, ignorability given the
coarsening $Q_n(Z)$; no claim of the latter is made or needed, since every
nuisance appearing in Section~\ref{sec:grid-inference} is compared against
the continuum truth $\ez^\circ,\mz^\circ$, never against a discretised
identifying assumption. Every background result of
Section~\ref{sec:background-results} that Section~\ref{sec:grid-inference}
invokes --- Proposition~\ref{prop:bilinear}, Lemma~\ref{lem:biasbound},
Lemma~\ref{lem:skeleton}, Lemma~\ref{lem:clip}, Lemma~\ref{lem:normequiv},
Lemma~\ref{lem:tangent} --- is instantiated at these continuum objects; each
holds verbatim under the $Z$-indexed conditions of this paragraph in place
of the $X$-indexed Assumptions~\ref{ass:causal}--\ref{ass:moments}, since
none of their proofs uses any property of the covariate space beyond what
this paragraph supplies (Lemma~\ref{lem:tangent} is already stated for an
arbitrary measurable covariate space; Proposition~\ref{prop:bilinear},
Lemma~\ref{lem:biasbound}, Lemma~\ref{lem:skeleton}, Lemma~\ref{lem:clip},
and Lemma~\ref{lem:normequiv} likewise never use finiteness or any other
structural property of $\cX$). What licenses substituting $Z$ for $X$ in
their proofs specifically is that $X_n=Q_n(Z)$ makes every function of
$X_n$ a function of $Z$, i.e.\ $\sigma(X_n)\subseteq\sigma(Z)$: each of
these results is an identity or bound involving nuisances and fitted
functions that are, in every instance below, measurable with respect to
whichever $\sigma$-field the result is stated over, and that measurability
transfers unchanged along the coarser-to-finer inclusion.

For each $j\in\{e,\mu\}$ there is a continuum axis-aligned tree $\tau_j^\circ$
with at most $\bar L$ leaves such that $\ez^\circ$ is constant on every leaf
of $\tau_e^\circ$ and $\mz^\circ$ is constant on every leaf of $\tau_\mu^\circ$.

The regime $d=d_n\to\infty$ is explicitly excluded: neither the overlap
bound nor the existence of a stable continuum limit for $\cX_n$ is
guaranteed to survive that regime; only mesh refinement at fixed $d$ is
considered.
\end{assumption}

The overlap bound is inherited automatically under refinement: since
$1-\ezn(x)=\E\{1-\ez^\circ(Z)\mid X_n=x\}\ge c$, cell averaging can only move
the conditional propensity into the convex hull of its within-cell values
and cannot weaken their common upper bound $1-c$.

Throughout, $\|\cdot\|_e,\|\cdot\|_\mu$ are the continuum-weighted norms of
Section~\ref{sec:setup}, $\|f\|_e^2=\E\{f(Z)^2\}$ and
$\|f\|_\mu^2=\E[\{1-\ez^\circ(Z)\}f(Z)^2]$, applied to a function of $X_n$ by
composition with $Q_n$; likewise $\nu$ denotes the weight with
$d\nu=\{1-\ez^\circ\}dP$. Define the grid-alignment errors, measured
directly against the \emph{continuum} truth,
\begin{equation}\label{eq:mesh-delta}
  \delta^*_{e,n}
  :=\min_{\tau\in\cT_{\bar L,n}}\|\Pi_\tau\ezn-\ez^\circ\|_e,
  \qquad
  \delta^*_{\mu,n}
  :=\min_{\tau\in\cT_{\bar L,n}}\|\Pi^\nu_\tau\mzn-\mz^\circ\|_\mu.
\end{equation}
For every $\tau\in\cT_{\bar L,n}$, every leaf of $\tau$ is a union of $X_n$
cells, hence a $\sigma(X_n)$-measurable event; by the tower property,
$\Pi_\tau\ezn=\Pi_\tau\ez^\circ$ and $\Pi^\nu_\tau\mzn=\Pi^\nu_\tau\mz^\circ$
as functions (the leaf-conditional mean of the once-projected $\ezn$
coincides with the leaf-conditional mean of $\ez^\circ$ itself, taken over
the same leaf). Consequently $\delta^*_{j,n}$ is exactly the best
$\le\bar L$-leaf, grid-restricted approximation error of the continuum
nuisance $\gamma_{0,j}^\circ$ itself, with no residual discretisation gap:
\[
  \delta^*_{j,n}=\min_{\tau\in\cT_{\bar L,n}}
  \|\Pi_\tau\gamma_{0,j}^\circ-\gamma_{0,j}^\circ\|_j .
\]

\begin{lemma}[Approximation under grid misalignment]\label{lem:mesh-approx}
Under Assumption~\ref{ass:mesh}, for $j\in\{e,\mu\}$,
\[
  \delta^*_{j,n}=O(\sqrt{h_n}).
\]
This rate concerns only the error caused by moving the finitely many
continuum split points of $\tau_j^\circ$ to the analyst's grid. It does not
assert that the best continuum approximation error at fixed $\bar L$
vanishes for a continuum nuisance that is not itself a tree step function
with at most $\bar L$ leaves.
\end{lemma}

\begin{proof}
\emph{Proof strategy.} Round each continuum split to a neighbouring grid
boundary without changing the tree topology or increasing its leaf count.
The rounded and continuum trees disagree only in finitely many slabs of
thickness at most $h_n$. Condition~\eqref{eq:mesh-slab} bounds their total
probability by $O(h_n)$; boundedness of the nuisances then gives squared
error $O(h_n)$ and norm error $O(\sqrt{h_n})$.

Fix $j\in\{e,\mu\}$ and write $\tau^\circ=\tau_j^\circ$. A binary tree with
at most $\bar L$ leaves has at most $\bar L-1$ internal splits. If its $s$th
split is at $z_{r_s}=t_s$, choose a grid boundary $t_{s,n}$ satisfying
$\abs{t_{s,n}-t_s}\le h_n$. Replacing every $t_s$ by $t_{s,n}$ yields a
grid-restricted tree $\tau_n^\circ\in\cT_{\bar L,n}$ with the same topology
after deletion of any empty leaves.

Let $B_n$ be the set on which $\tau^\circ$ and $\tau_n^\circ$ assign
different leaves. A change of assignment can occur only if at least one
rounded split is crossed, hence
\[
  B_n\subseteq\bigcup_{s=1}^{\card{\tau^\circ}-1}
  \{z:\abs{z_{r_s}-t_s}\le h_n\}.
\]
The union bound and~\eqref{eq:mesh-slab} give
\begin{align}
  \PP(Z\in B_n)
  &\le\sum_{s=1}^{\card{\tau^\circ}-1}\PP\{\abs{Z_{r_s}-t_s}\le h_n\}\\
  &\le(\bar L-1)C_Xh_n=O(h_n).
  \label{eq:mesh-boundary-mass}
\end{align}

For the propensity, define $g^\circ_e$ by assigning to each leaf of
$\tau_n^\circ$ the value carried by the corresponding leaf of $\tau_e^\circ$.
Away from $B_n$, the value of $g^\circ_e$ agrees with $\ez^\circ$. Because
$0\le\ez^\circ\le1$,
\[
  \{g^\circ_e(Z)-\ez^\circ(Z)\}^2\le\mathbf1\{Z\in B_n\}.
\]
Taking expectations,
$\E\{g^\circ_e(Z)-\ez^\circ(Z)\}^{2}\le\PP(Z\in B_n)=O(h_n)$.
Since $g^\circ_e$ is itself constant on the leaves
of $\tau_n^\circ$, and $\Pi_{\tau_n^\circ}\ez^\circ$ is the $\|\cdot\|_e$-best
approximation to $\ez^\circ$ among such functions,
\[
  \|\Pi_{\tau_n^\circ}\ez^\circ-\ez^\circ\|_e^2
  \le\E\{g^\circ_e(Z)-\ez^\circ(Z)\}^{2}\le O(h_n).
\]
Every leaf of $\tau_n^\circ$ is a union of $X_n$ cells, hence
$\sigma(X_n)$-measurable, so by the tower property
$\Pi_{\tau_n^\circ}\ezn=\Pi_{\tau_n^\circ}\ez^\circ$. By~\eqref{eq:mesh-delta},
\[
  (\delta^*_{e,n})^2
  \le\|\Pi_{\tau_n^\circ}\ezn-\ez^\circ\|_e^{\,2}\le O(h_n).
\]

For the outcome nuisance, $\abs{\mz^\circ}\le B_Y$. The same construction
gives a candidate $g^\circ_\mu$ satisfying
$\{g^\circ_\mu(Z)-\mz^\circ(Z)\}^2\le4B_Y^2\mathbf1\{Z\in B_n\}$. Taking
expectations weighted by $\{1-\ez^\circ\}\le1$,
\[
  \E\bigl[\{1-\ez^\circ(Z)\}\{g^\circ_\mu(Z)-\mz^\circ(Z)\}^{2}\bigr]
  \le4B_Y^2\PP(Z\in B_n)=O(h_n).
\]
By the same argument, $\Pi_{\tau_n^\circ}^\nu\mz^\circ$ is the
$\|\cdot\|_\mu$-best approximation to $\mz^\circ$ among functions constant
on the leaves of $\tau_n^\circ$, and $g^\circ_\mu$ is one such function, so
$\|\Pi_{\tau_n^\circ}^\nu\mz^\circ-\mz^\circ\|_\mu^2\le4B_Y^2\PP(Z\in B_n)
=O(h_n)$; the tower property gives
$\Pi_{\tau_n^\circ}^\nu\mzn=\Pi_{\tau_n^\circ}^\nu\mz^\circ$, and hence
$(\delta^*_{\mu,n})^2\le O(h_n)$. Taking square roots proves both claims.

The rounding construction uses the vocabulary of free-knot approximation
\citep{deboor1973,barrowsmith1978}, though not their result type: those
papers give asymptotic rates in the number of knots under smoothness of the
target function, whereas here a \emph{fixed} number of knots is perturbed
by at most $h_n$, with no smoothness assumed. Its use for histogram and
data-dependent-partition approximations follows the same geometric
principle as \citet{nobel1996} and \citet{lugosinobel1996}. The only
imported idea is the replacement of free split points by nearby admissible
knots; the probability and norm bounds above are derived directly under
Assumption~\ref{ass:mesh}.
\end{proof}

\begin{remark}[Post-proof audit of Lemma~\ref{lem:mesh-approx}]
\emph{Assumption check.} Fixed $\bar L$ bounds the number of slabs;
\eqref{eq:mesh-slab} controls their mass; bounded nuisances control the
error inside them.
\emph{Dependency check.} Only the union bound, conditional Jensen, and
projection optimality are used in-line; the cited approximation literature
supplies context rather than an unproved logical step.
\emph{Rate verification.} Boundary mass is $O(h_n)$, hence squared error is
$O(h_n)$ and norm error is $O(\sqrt{h_n})$.
\emph{Edge-case scan.} The rate can fail for atoms on true split boundaries
or distributions violating~\eqref{eq:mesh-slab}.
\emph{Tightness.} For a non-zero jump and a bounded density positive near
the split, $O(\sqrt{h_n})$ is generally sharp in $L_2$.
\emph{Failure localization.} The vulnerable step is
\eqref{eq:mesh-boundary-mass}; a weaker boundary-mass modulus would yield
the corresponding weaker norm rate.
\end{remark}

\begin{remark}[Scope of the approximation result]
The displayed model never exceeds $\bar L$ leaves, and exact sufficiency is
required only of the continuum-level nuisances: true split points may lie
anywhere, not on the analyst's discretisation grid, and
Lemma~\ref{lem:mesh-approx} shows the resulting grid-alignment error
vanishes at rate $O(\sqrt{h_n})$ as the grid refines, using only a
boundary-mass union bound --- no free-knot rate, no VC dimension, no
chaining, no localisation. This result does not claim that the fixed-budget
continuum approximation bias vanishes when the continuum truth is not
itself representable by $\le\bar L$ leaves; it claims only that the
specific bias induced by grid misalignment can be driven to zero by
refining the grid, without ever increasing $\bar L$. This is a
population-level, approximation-theoretic statement; it makes no claim
about inference. Section~\ref{sec:grid-inference} takes up that question.
\end{remark}

% ===========================================================================
\section{Inference under grid refinement}\label{sec:grid-inference}

\subsection{Why exact selection consistency cannot survive, and what
replaces it}\label{sec:grid-inference-why}

The fixed-grid theory's route to structure-selection consistency
(companion document's selection lemma) fixes a tolerance $\delta<\Delta_j/2$
against the population risk margin $\Delta_j$ --- the excess risk of the
best insufficient candidate over the best candidate --- which is a strictly
positive constant under a fixed grid. Under Assumption~\ref{ass:mesh}, the
excluded class $\cT_{\bar L,n}\setminus\cS_{j,n}$ comes to contain trees one
grid boundary short of a true continuum jump, so the corresponding margin
degrades to $\Delta_{j,n}=O(h_n)\to0$. Rescuing the fixed-grid argument by a
uniformity assumption on $\Delta_{j,n}$ would require
$\bar L\log(d/h_n)=o(nh_n^2)$; at the target rate $h_n=o(n^{-1/2})$ the
right side tends to zero while the left side diverges like $\bar L\log n$.
No choice of constants repairs this: exact selection consistency is not
merely harder to establish under refinement, it is false in exactly the
regime this document targets. What follows replaces it with a genuine
finite-class penalized-ERM oracle inequality --- an excess-risk bound, not a
membership statement --- built from a localized Bernstein/curvature
argument rather than any margin comparison.

A second, independent obstruction is worth stating plainly before
proceeding. Structural sparsity (the companion document's exact-sufficiency
assumption) requires $\cS_{e,n},\cS_{\mu,n}\ne\emptyset$ at each row $n$.
Under grid refinement this fails at \emph{every finite $n$}, not merely
asymptotically: the discretized truth $\ezn,\mzn$ is generically not exactly
constant on any grid-restricted tree's leaves, since $\delta^*_{j,n}>0$
whenever the true continuum split does not happen to land on the row-$n$
grid. Consequently every part of the companion document's efficiency
machinery that lists exact sparsity as a hypothesis is simply inapplicable
here, at any row, and cannot be salvaged by any row-wise reading of its
proofs. What remains available without sparsity is the companion
document's cross-fitted theorem, anchor interval, and width corollary
(Section~\ref{sec:background-results} above); the results below either
build on those directly or independently establish what the sparsity-based
machinery cannot.

Put
\[
  s_n:=1+\log\card{\cT_{\bar L,n}},
  \qquad
  r_n:=\sqrt{\frac{s_n}{n}}.
\]

\begin{assumption}[Uniform refining-grid inferential conditions]
\label{ass:grid-inference}
In addition to Assumption~\ref{ass:mesh}, the following conditions hold.

\begin{enumerate}[label=(U\arabic*),leftmargin=*]

\item\label{u:entropy}
There is a constant $C_T<\infty$, independent of $n$, such that
\[
  \log\card{\cT_{\bar L,n}}
  \le C_T\bar L\log(dJ_{n,\max}),
  \qquad
  s_n=o(n).
\]
For the root-$n$ conclusions below we impose the stronger condition
\[
  s_n=o(\sqrt n).
  \tag{U1$^+$}\label{eq:u1plus}
\]
When $J_{n,\max}\asymp h_n^{-1}$ and $h_n$ is polynomial in $n$,
$s_n=O\{\bar L\log(d/h_n)\}=O(\log n)$, so both conditions hold.

\item\label{u:comparator}
For each $j\in\{e,\mu\}$, let $\tau^{\mathrm{comp}}_{j,n}$ be the
grid-rounded tree constructed in the proof of
Lemma~\ref{lem:mesh-approx}. After deleting $P_Z$-null continuum leaves,
there is a constant $q_{\min}>0$, independent of $n$, such that, for all
sufficiently large $n$,
\[
  \min_{\ell\in\tau^{\mathrm{comp}}_{j,n}}P\{X_n\in\ell\}\ge q_{\min}.
\]
Moreover $m_n/n\to0$, so these comparator trees belong to the feasible sets
in~\eqref{eq:feasible} with probability tending to one.

\item\label{u:bernstein}
The squared-loss fits are uniformly bounded: propensity fits lie in $[0,1]$
before and after clipping, and outcome fits lie in $[-B_Y,B_Y]$. The
penalty satisfies $\bar L\lambda_n=O(r_n^2)$. For each nuisance and each
training sample of order $n$, the finite union of at-most-$\bar
L$-dimensional leaf-constant models satisfies the localized quadratic-loss
bound
\begin{equation}\label{eq:grid-bernstein}
 \sup_{f\in\mathcal F_{j,n}}
 \frac{\left|(P_n-P)\{\ell_j(f)-\ell_j(\gamma_{0,j}^\circ)\}\right|}
      {\|f-\gamma_{0,j}^\circ\|_j r_n+r_n^2}
 =O_p(1),
\end{equation}
where $\ell_e(f)=(A-f(X_n))^2$, $\ell_\mu(f)=(1-A)\{Y-f(X_n)\}^2$, and
$\mathcal F_{j,n}$ is the union, over $\tau\in\cT_{\bar L,n}$, of the
bounded functions constant on the leaves of $\tau$.

\item\label{u:approx}
The grid-alignment errors satisfy $\delta^*_{j,n}=O(\sqrt{h_n})$ for
$j\in\{e,\mu\}$, as established by Lemma~\ref{lem:mesh-approx}.

\end{enumerate}
\end{assumption}

\paragraph{Why (U2) is derived rather than postulated.} Let $\ell^\circ$ be
a positive-mass leaf of one of the fixed continuum trees $\tau_j^\circ$, and
let $\ell_n$ be its rounded counterpart. The proof of
Lemma~\ref{lem:mesh-approx} constructs a disagreement set $B_n$ satisfying
$P_Z(B_n)\le(\bar L-1)C_Xh_n$. Since $\ell_n\triangle\ell^\circ\subseteq
B_n$,
\[
  |P_Z(\ell_n)-P_Z(\ell^\circ)|
  \le P_Z(\ell_n\triangle\ell^\circ)
  \le(\bar L-1)C_Xh_n.
\]
There are finitely many positive-mass continuum leaves, so
$q^\circ_{\min}:=\min_{j\in\{e,\mu\}}\min_{\ell^\circ\in\tau_j^\circ:
P_Z(\ell^\circ)>0}P_Z(\ell^\circ)>0$, and for all sufficiently large $n$,
$P_Z(\ell_n)\ge q^\circ_{\min}/2$. Thus (U2) holds with
$q_{\min}=q^\circ_{\min}/2$. By $1-\ez^\circ\ge c$, every such leaf has
control mass at least $cq_{\min}$. Finally, $m_n/n\to0$ and the weak law of
large numbers imply that each comparator leaf eventually contains at least
$m_n$ observations (and at least $m_n$ controls for the outcome fit) with
probability tending to one.

\subsection{The oracle inequality}\label{sec:grid-oracle}

\begin{lemma}[Refining-grid oracle inequality]
\label{lem:grid-oracle}
Suppose Assumptions~\ref{ass:iid}, \ref{ass:causal},
\ref{ass:construct}, \ref{ass:global}, \ref{ass:mesh}, and
\ref{ass:grid-inference} hold. Let $\widehat\gamma_{j,n}$ be the penalized
empirical-risk minimizer over the row-$n$ feasible class. Then, for each
$j\in\{e,\mu\}$,
\begin{equation}\label{eq:grid-oracle}
  \|\widehat\gamma_{j,n}-\gamma_{0,j}^\circ\|_j
  =O_p\!\left(\sqrt{h_n}+r_n\right)
  =O_p\bigl(\delta^*_{j,n}+r_n\bigr),
\end{equation}
the second equality because Lemma~\ref{lem:mesh-approx}, under the
corrected definition~\eqref{eq:mesh-delta}, gives $\delta^*_{j,n}=O(\sqrt{h_n})$
via exactly the comparator $\tau^{\mathrm{comp}}_{j,n}$ used in the proof
below; no separate leaf-mass property beyond (U2) is required for this
equivalence. After clipping,
\begin{equation}\label{eq:grid-weight-oracle}
  D_{w,n}=O_p(\sqrt{h_n}+r_n).
\end{equation}
\end{lemma}

\begin{proof}
\emph{Proof strategy.} Use the grid-rounded tree as an actually feasible
comparator. Apply the penalized ERM basic inequality, control its
empirical-process term by the localized Bernstein bound
\eqref{eq:grid-bernstein}, and absorb the resulting linear term into the
quadratic population excess risk. This argument compares risks and never
compares a shrinking population margin.

Fix $j$. Write $g_j=\gamma_{0,j}^\circ$, $\widehat f=\widehat\gamma_{j,n}$,
and let $f^{\mathrm{comp}}_{j,n}$ be the population leaf projection on
$\tau^{\mathrm{comp}}_{j,n}$. By (U2), the comparator is feasible with
probability tending to one. On that event, optimality in~\eqref{eq:select}
gives
\begin{align}
 P_n\ell_j(\widehat f)+\lambda_n\card{\widehat\tau_j}
 &\le
 P_n\ell_j(f^{\mathrm{comp}}_{j,n})
   +\lambda_n\card{\tau^{\mathrm{comp}}_{j,n}},\\
 P_n\{\ell_j(\widehat f)-\ell_j(g_j)\}
 &\le
 P_n\{\ell_j(f^{\mathrm{comp}}_{j,n})-\ell_j(g_j)\}
   +\bar L\lambda_n.
 \label{eq:grid-basic}
\end{align}

For squared loss and a conditional-mean truth, the population excess-risk
identity is exact. For the propensity,
\[
 P\{\ell_e(f)-\ell_e(\ez^\circ)\}
 =E\{f(X_n)-\ez^\circ(Z)\}^2=\|f-\ez^\circ\|_e^2,
\]
since expanding the squares gives
$(A-f)^2-(A-\ez^\circ)^2=(f-\ez^\circ)^2-2(A-\ez^\circ)(f-\ez^\circ)$, and
the second term has expectation zero because $E(A-\ez^\circ\mid Z)=0$ and
$f(X_n)$ is $\sigma(Z)$-measurable. Similarly,
\[
 P\{\ell_\mu(f)-\ell_\mu(\mz^\circ)\}
 =E[\{1-\ez^\circ(Z)\}\{f(X_n)-\mz^\circ(Z)\}^2]
 =\|f-\mz^\circ\|_\mu^2,
\]
because $E[(1-A)\{Y-\mz^\circ(Z)\}\mid Z]=0$.

Let $d=\|\widehat f-g_j\|_j$, $d_c=\|f^{\mathrm{comp}}_{j,n}-g_j\|_j$.
Adding and subtracting population risks in~\eqref{eq:grid-basic}, and
applying \eqref{eq:grid-bernstein} to $\widehat f$ and the comparator,
yields
\[
 d^2
 \le d_c^2+O_p(dr_n+r_n^2)+O_p(d_cr_n+r_n^2)
       +\bar L\lambda_n.
\]
By (U3), $\bar L\lambda_n=O(r_n^2)$. Young's inequality,
$ab\le a^2/4+b^2$, gives $O_p(dr_n)\le\frac14d^2+O_p(r_n^2)$ and
$O_p(d_cr_n)\le\frac14d_c^2+O_p(r_n^2)$. Consequently
$\frac34d^2\le\frac54d_c^2+O_p(r_n^2)$, so
\[
 d^2=O_p(d_c^2+r_n^2),
 \qquad
 d=O_p(d_c+r_n).
\]
The comparator construction in Lemma~\ref{lem:mesh-approx} gives
$d_c=O(\sqrt{h_n})$, proving the first equality in~\eqref{eq:grid-oracle}.
The second equality is immediate since, by Lemma~\ref{lem:mesh-approx} and
the corrected definition~\eqref{eq:mesh-delta}, $d_c$ and $\delta^*_{j,n}$
are controlled by the identical bound arising from the identical
comparator: no separate "sharper" case or additional leaf-mass property is
needed.

Finally, on $[0,1-c]$ the map $u\mapsto u/(1-u)$ has derivative at most
$c^{-2}$. The mean-value theorem, Lemma~\ref{lem:clip}, and
Lemma~\ref{lem:normequiv} therefore give
\[
 D_{w,n}
 \le c^{-2}\|\widehat e_n-\ez^\circ\|_\mu
 \le c^{-2}\|\widehat e_n-\ez^\circ\|_e
 =O_p(\sqrt{h_n}+r_n).
\]
\end{proof}

\paragraph{Post-proof audit.}
\emph{Assumptions:} boundedness supplies the Bernstein envelope; (U1)
controls the number of partition models; (U2) makes the comparator
feasible; (U3) controls both local stochastic fluctuations and the
penalty; and (U4) supplies its approximation rate.
\emph{Dependencies:} only the squared-loss excess-risk identity,
\eqref{eq:grid-bernstein}, Young's inequality, Lemma~\ref{lem:mesh-approx},
Lemma~\ref{lem:clip}, and Lemma~\ref{lem:normequiv} are used.
\emph{Rates:} the quadratic inequality yields norm error
$O_p(\sqrt{h_n}+r_n)$, not $O_p(\sqrt{h_n}+n^{-1/2})$ unless $s_n=O(1)$.
\emph{Edge cases:} the result can fail if the comparator loses leaf mass,
the outcome is unbounded without a replacement tail condition, or the
penalty dominates $r_n^2$.
\emph{Tightness:} the $\sqrt{s_n/n}$ model-selection term is generally
unavoidable over a growing finite class.
\emph{Failure localization:} the substantive condition is
\eqref{eq:grid-bernstein}; it is the Bernstein/curvature replacement for the
invalid shrinking-margin argument. No object is treated as invariant across
$n$: the grid, class, comparator, projections, and fitted functions all
carry an explicit row index.

\subsection{Result A: cross-fitted inference (a true corollary)}\label{sec:result-a}

\begin{corollary}[Cross-fitted inference on a refining grid]
\label{cor:grid-crossfit}
Suppose the conditions of Lemma~\ref{lem:grid-oracle} hold for every
fold-wise training sample, with fixed $K$. Suppose additionally that
\begin{equation}
  h_n=o(n^{-1/2}),
  \qquad
  s_n=o(\sqrt n).
  \label{eq:grid-rootn}
\end{equation}
Let $\thetatil_n$ be the cross-fitted estimator of Theorem~\ref{thm:crossfit},
constructed from the row-$n$ grids. Then
\[
  \sqrt n(\thetatil_n-\theta_0)\rightsquigarrow N(0,V).
\]
\end{corollary}

\begin{proof}
\emph{Proof strategy.} Verify Assumption~\ref{ass:rate} using
Lemma~\ref{lem:grid-oracle}, and then apply Theorem~\ref{thm:crossfit}
verbatim. Cross-fitting removes any need for uniform empirical-process
control at the score-evaluation stage.

For each fixed fold, its training sample has size $n(1-1/K)+O(1)\asymp n$,
so Lemma~\ref{lem:grid-oracle} gives
$\|\widehat e_n-\ez^\circ\|_e=O_p(\sqrt{h_n}+r_n)$ and
$\|\widehat\mu_n-\mz^\circ\|_\mu=O_p(\sqrt{h_n}+r_n)$. Their product is
\[
 \|\widehat e_n-\ez^\circ\|_e\|\widehat\mu_n-\mz^\circ\|_\mu
 =O_p\{(\sqrt{h_n}+r_n)^2\}
 =O_p\left(h_n+2\sqrt{h_n}\,r_n+r_n^2\right).
\]
Multiplying each deterministic rate by $\sqrt n$ gives $\sqrt n\,h_n\to0$
and $\sqrt n\,r_n^2=s_n/\sqrt n\to0$ by~\eqref{eq:grid-rootn}, and
\[
 \sqrt n\,\sqrt{h_n}\,r_n
 =\sqrt{h_ns_n}
 =\sqrt{(\sqrt n h_n)(s_n/\sqrt n)}
 \longrightarrow0.
\]
Hence $\|\widehat e_n-\ez^\circ\|_e\|\widehat\mu_n-\mz^\circ\|_\mu
=o_p(n^{-1/2})$, which is precisely Assumption~\ref{ass:rate}, with the
covariate in the score understood to be $Z$ and the fitted nuisances
restricted to functions of $X_n=Q_n(Z)$.

The remaining hypotheses of Theorem~\ref{thm:crossfit} --
Assumptions~\ref{ass:iid}, \ref{ass:causal}, and~\ref{ass:moments} --
are in force here in their $Z$-indexed form, established in
Assumption~\ref{ass:mesh}'s causal-identification paragraph, not the
$X$-indexed form of Section~\ref{sec:background-assumptions};
Assumptions~\ref{ass:construct} and~\ref{ass:global} restrict the
estimator's construction and are unaffected by which covariate identifies
$\theta_0$, and are included above as stated. Its proof is grid-free after
conditioning on each training fit: the held-out empirical-process term is a
centred average independent of that fit. Theorem~\ref{thm:crossfit}
therefore yields $\sqrt n(\thetatil_n-\theta_0)\rightsquigarrow N(0,V)$.
\end{proof}

\paragraph{Post-proof audit.}
\emph{Assumptions:} the stronger entropy requirement $s_n=o(\sqrt n)$ is
necessary because the oracle norm error contains $\sqrt{s_n/n}$.
\emph{Dependencies:} Lemma~\ref{lem:grid-oracle} verifies
Assumption~\ref{ass:rate}; Theorem~\ref{thm:crossfit} then applies
directly.
\emph{Rates:} all three terms in $(\sqrt{h_n}+r_n)^2$ are explicitly
$o(n^{-1/2})$.
\emph{Edge cases:} U1 alone, $s_n=o(n)$, is insufficient; it permits
$s_n/\sqrt n\nrightarrow0$.
\emph{Tightness:} the mesh condition is sufficient and matches the product
of two $O(\sqrt{h_n})$ approximation errors.
\emph{Failure localization:} failure can occur at the DML product-rate
step, not in Theorem~\ref{thm:crossfit}. Nothing is assumed literally
invariant across rows except the fixed constants $d,\bar L,c,B_Y,C_T$, and
$q_{\min}$.

\subsection{Result A$'$: anchor-interval width (a rate, not an efficiency,
statement)}\label{sec:result-aprime}

\begin{corollary}[Anchor-interval width under grid refinement]
\label{cor:grid-width}
Suppose the conditions of Corollary~\ref{cor:grid-crossfit} hold and the
cross-fitted variance estimator satisfies $\widetilde\sigma_n^2\to_pV$.
Form the anchor interval of Theorem~\ref{thm:anchor} using $\thetatil_n$.
Then
\[
 \liminf_{n\to\infty}P(\theta_0\in C_n)\ge1-\alpha
\]
and
\begin{equation}\label{eq:grid-width}
 \operatorname{width}(C_n)
 =O_p\!\left(
   \max\left\{
     \sqrt{\frac{\bar L\log(dJ_{n,\max})}{n}},
     D_{w,n}D_{\mu,n}
   \right\}
 \right).
\end{equation}
Under Lemma~\ref{lem:grid-oracle},
$D_{w,n}D_{\mu,n}=O_p\left(h_n+\sqrt{h_n}\,r_n+r_n^2\right)$. \textbf{This is
a width-rate statement, not an efficiency statement.}
\end{corollary}

\begin{proof}
\emph{Proof strategy.} Use Theorem~\ref{thm:anchor} for validity. Repeat
the decomposition in the proof of Corollary~\ref{cor:width}, replacing its
fixed-class union bound by the growing-class rate $r_n$.

Theorem~\ref{thm:anchor} applies because Corollary~\ref{cor:grid-crossfit}
supplies a valid asymptotically normal anchor and
$\widetilde\sigma_n^2\to_pV>0$. Validity therefore follows without any
condition on the full-sample tree estimator.

For width, write $\widehat\delta_n=(\widehat\theta_n-\theta_0)
-(\thetatil_n-\theta_0)$. By Corollary~\ref{cor:grid-crossfit},
$\thetatil_n-\theta_0=O_p(n^{-1/2})$. Proposition~\ref{prop:bilinear} and
Lemma~\ref{lem:biasbound} give
$|\theta(\widehat\eta_n)-\theta_0|\le\pi^{-1}D_{w,n}D_{\mu,n}$. The
fixed-class $O_p(n^{-1/2})$ plug-in fluctuation used in
Corollary~\ref{cor:width} becomes $O_p(r_n)$ under (U1) and
\eqref{eq:grid-bernstein}. Therefore
$|\widehat\theta_n-\theta_0|=O_p\{D_{w,n}D_{\mu,n}+r_n\}$, and hence
$|\widehat\delta_n|=O_p\{D_{w,n}D_{\mu,n}+r_n\}$. Since
$\operatorname{width}(C_n)=2\{|\widehat\delta_n|
+z_{1-\alpha/2}\widetilde\sigma_n/\sqrt n\}$ and $r_n\ge n^{-1/2}$ by
definition of $s_n$, this proves
$\operatorname{width}(C_n)=O_p\{\max(r_n,D_{w,n}D_{\mu,n})\}$. Finally,
Lemma~\ref{lem:grid-oracle} gives $D_{w,n}D_{\mu,n}=O_p\{(\sqrt{h_n}+r_n)^2\}
=O_p(h_n+\sqrt{h_n}\,r_n+r_n^2)$. Substitution and (U1) yield
\eqref{eq:grid-width}.
\end{proof}

\paragraph{Post-proof audit.}
\emph{Assumptions:} anchor validity and width control are logically
separate; only the latter uses the oracle rates.
\emph{Dependencies:} Theorem~\ref{thm:anchor}, Corollary~\ref{cor:grid-crossfit},
Proposition~\ref{prop:bilinear}, and Lemma~\ref{lem:biasbound} are used.
\emph{Rates:} the fixed-class $n^{-1/2}$ term is replaced by
$r_n=\sqrt{s_n/n}$.
\emph{Edge cases:} coverage survives even when the width fails to contract,
provided the anchor remains valid.
\emph{Tightness:} no matching lower bound for the width is asserted.
\emph{Failure localization:} \textbf{this corollary does not establish
$\sqrt n(\widehat\theta_n-\thetatil_n)=o_p(1)$ and therefore does not
establish efficiency of the displayed-tree estimator.} Every
class-dependent rate explicitly retains its $n$ index.

\subsection{Result B: full-sample inference (an independent theorem)}\label{sec:result-b}

\begin{theorem}[Full-sample inference under vanishing grid misfit]
\label{thm:grid-fullsample}
Suppose the conditions of Corollary~\ref{cor:grid-crossfit} hold. Suppose,
in addition, that the following three conditions hold.
Hypotheses~\ref{h:wclass}--\ref{h:aclass} are marginal instances of the
localized-concentration technique behind (U3), not a single joint
hypothesis over pairs of trees; Hypothesis~\ref{h:refit} additionally
carries a new primitive rate condition, discussed in the post-proof audit
below.

\begin{enumerate}[label=(H\arabic*),leftmargin=*]
\item\label{h:wclass} Writing $\widehat w_\tau$ for the propensity
leaf-refit fit on candidate partition $\tau$, uniformly over
$\tau\in\cT_{\bar L,n}$,
\[
  \frac{\Bigl|n^{-1}\sum_{i:A_i=0}\{\widehat w_\tau(Z_i)-w_0^\circ(Z_i)\}^2
  -P[(1-A)(\widehat w_\tau-w_0^\circ)^2]\Bigr|}
  {\|\widehat w_\tau-w_0^\circ\|_\mu\,r_n+r_n^2}=O_p(1).
\]
(This is the relative/ratio form used throughout (U3); an \emph{absolute}
bound of $O_p(D_{w,n}^2+r_n^2)$ cannot hold uniformly over $\tau$ -- take
$\tau$ the root-only tree, for which the empirical squared error
$n^{-1}\sum_{i:A_i=0}\{\widehat w_\tau(Z_i)-w_0^\circ(Z_i)\}^2$ itself is
$\Theta_p(1)$ while $D_{w,n}^2+r_n^2\to0$. Evaluated at the selected
$\widehat\tau_e$, this gives $n^{-1}\sum_{i:A_i=0}\{\widehat w_{\widehat\tau_e}(Z_i)
-w_0^\circ(Z_i)\}^2\le\|\widehat w_{\widehat\tau_e}-w_0^\circ\|_\mu^2
+O_p(D_{w,n}r_n+r_n^2)=O_p(D_{w,n}^2+r_n^2)$ by Young's inequality, since
$\|\widehat w_{\widehat\tau_e}-w_0^\circ\|_\mu=D_{w,n}$ exactly -- which is
all Step 2 below uses.)
\item\label{h:aclass} Writing $a_{\mu,\tau}:=\Pi^\nu_\tau\mzn-\mz^\circ$
for the (deterministic) population leaf-projection residual and
$h=w_0^\circ(1-A)-A$, uniformly over the deterministic family
$\{a_{\mu,\tau}:\tau\in\cT_{\bar L,n}\}$,
\[
  |P_n[h\,a_{\mu,\tau}]|=O_p(\|a_{\mu,\tau}\|_\mu r_n+r_n^2)
  \quad\text{and}\quad
  \bigl|P_n[(1-A)a_{\mu,\tau}^2]-P[(1-A)a_{\mu,\tau}^2]\bigr|
  =O_p(\|a_{\mu,\tau}\|_\mu r_n+r_n^2).
\]
(The second display is stated in the same relative form as the first, for
the same reason as Hypothesis~\ref{h:wclass}: an absolute $O_p(r_n^2)$ bound
is false at the root-only tree, and indeed is falsified \emph{by}
(U1$^+$), since there the true deviation is $\Theta_p(n^{-1/2})\gg r_n^2$.)
\item\label{h:refit} Writing $v_{\mu,\tau}:=\widehat\mu_\tau-\Pi^\nu_\tau\mzn$
for the leaf-refit residual, uniformly over $\tau\in\cT_{\bar L,n}$,
\[
  \sum_{\ell\in\tau}\bigl|v_{\ell,\tau}P_n[h\mathbf1_\ell]\bigr|=O_p(r_n^2),
\]
which we take together with the additional rate condition $s_n=o(m_n)$ on
the per-leaf sample floor $m_n$ of Assumption~\ref{ass:construct}(d); this
condition is new relative to Assumption~\ref{ass:grid-inference} and is
needed only for this hypothesis.
\end{enumerate}

Then the full-sample estimator $\widehat\theta_n$ formed from the two
selected row-$n$ trees satisfies
\[
 \sqrt n(\widehat\theta_n-\theta_0)
 =\sqrt n\,P_n\psi_0+o_p(1)
 \rightsquigarrow N(0,V).
\]
Moreover, the plug-in estimator~\eqref{eq:vhat-def} satisfies
$\widehat V_n\to_pV$, and the corresponding Wald interval has pointwise
asymptotic coverage $1-\alpha$.

If, additionally, the interiority condition of Lemma~\ref{lem:tangent}
holds, the contiguity argument of Theorem~\ref{thm:regular} applies to the
displayed asymptotic-linear expansion; thus $\widehat\theta_n$ is regular
and locally asymptotically efficient at the fixed continuum law $P_0$.
\end{theorem}

\begin{equation}\label{eq:vhat-def}
  \widehat V_n
  =\widehat\pi_n^{-2}\,n^{-1}\sum_{i=1}^n
  \Bigl[A_i\bigl\{Y_i-\widehat\mu_n(X_{n,i})-\widehat\theta_n\bigr\}
   -(1-A_i)\widehat w_n(X_{n,i})\bigl\{Y_i-\widehat\mu_n(X_{n,i})\bigr\}\Bigr]^2 .
\end{equation}

\begin{proof}
\emph{Proof strategy.} Start from the exact skeleton identity of
Lemma~\ref{lem:skeleton}. Replace the sparsity-dependent proofs of the
companion document's leaf-value and cross terms by localized bounds in the
vanishing realised errors. The genuinely new step is uniformity over the
growing, data-selected partition class. After proving the two remainders
are $o_p(n^{-1/2})$, apply the ordinary central limit theorem. Finally
prove variance consistency directly from vanishing empirical $L_2$ error;
the fixed-misspecification failure of Remark~\ref{rem:honest} does not
apply because both discrepancies now vanish.

Let $a_{\mu,\tau}:=\Pi^\nu_\tau\mzn-\mz^\circ$ and
$v_{\mu,\tau}:=\widehat\mu_\tau-\Pi^\nu_\tau\mzn$, so
$\widehat\mu_\tau-\mz^\circ=a_{\mu,\tau}+v_{\mu,\tau}$ (matching the notation
of Hypotheses~\ref{h:aclass}--\ref{h:refit} above).

\emph{Step 1: the leaf-value term.} For fixed $\tau$, decompose
$U(\tau)=P_n[h\,a_{\mu,\tau}]+P_n[h\,v_{\mu,\tau}]$. Since $E(h\mid Z)=0$ by
Lemma~\ref{lem:skeleton}, $P[h\,a_{\mu,\tau}]=0$. Hypothesis~\ref{h:aclass}'s
first display gives, simultaneously for all $\tau$,
\begin{equation}\label{eq:grid-U-approx}
 |P_n[h\,a_{\mu,\tau}]|
 =O_p\{\|a_{\mu,\tau}\|_\mu r_n+r_n^2\}.
\end{equation}
For the leaf-refit component, write
$v_{\mu,\tau}=\sum_{\ell\in\tau}v_{\ell,\tau}\mathbf1_\ell$; then
$P_n[h\,v_{\mu,\tau}]=\sum_{\ell\in\tau}v_{\ell,\tau}P_n[h\mathbf1_\ell]$.
Hypothesis~\ref{h:refit} gives, uniformly over $\tau$,
\begin{equation}\label{eq:grid-U-refit}
 \sum_{\ell\in\tau}|v_{\ell,\tau}P_n[h\mathbf1_\ell]|=O_p(r_n^2).
\end{equation}
This is the growing-class analogue of the companion document's leaf-value
bound; it uses neither leaf-constancy of $\mz^\circ$ nor structural
sparsity. Evaluating \eqref{eq:grid-U-approx}--\eqref{eq:grid-U-refit} at
the selected partition and using Lemma~\ref{lem:grid-oracle},
$|U|=O_p\{D_{\mu,n}r_n+r_n^2\}$, so
\[
 \sqrt n\,|U|
 =O_p\{D_{\mu,n}\sqrt{s_n}+s_n/\sqrt n\}.
\]
Now $D_{\mu,n}\sqrt{s_n}=O_p\{\sqrt{h_ns_n}+s_n/\sqrt n\}=o_p(1)$ by
\eqref{eq:grid-rootn}. Hence $\sqrt n\,U=o_p(1)$.

\emph{Step 2: the cross term.} For each fixed pair $(\tau_e,\tau_\mu)$,
centre $\Delta w_\tau=\widehat w_{\tau_e}-w_0^\circ$ within the control
observations of each $\tau_\mu$ leaf:
$b_i=\Delta w_\tau(Z_i)-\overline{\Delta w}_{\ell_\mu(i)}$. The exact refit
identity gives
$T(\tau_e,\tau_\mu)=n^{-1}\sum_{i:A_i=0}b_i\{Y_i-\widehat\mu_{\tau_\mu}(X_{n,i})\}$.
Insert
$Y_i-\widehat\mu_{\tau_\mu}=\varepsilon_i-a_{\mu,\tau_\mu}(Z_i)-v_{\mu,\tau_\mu}(X_{n,i})$,
$\varepsilon_i=Y_i-\mz^\circ(Z_i)$. The term containing $v_{\mu,\tau_\mu}$
vanishes exactly because it is constant within each $\tau_\mu$ leaf and
$\sum_{\ell,A_i=0}b_i=0$. Therefore $T=T_\varepsilon-T_a$, where
$T_\varepsilon=n^{-1}\sum_{i:A_i=0}b_i\varepsilon_i$ and
$T_a=n^{-1}\sum_{i:A_i=0}b_ia_{\mu,\tau_\mu}(Z_i)$.

Fix a pair $(\tau_e,\tau_\mu)$. Conditionally on $(Z_{1:n},A_{1:n})$, the
coefficients $b_i$ are measurable, the $\varepsilon_i$ are independent and
mean zero, and $E(\varepsilon_i^2\mid Z_i,A_i=0)\le B_Y^2$. Hence
$\operatorname{Var}(T_\varepsilon\mid Z,A)\le B_Y^2n^{-2}\sum_{i:A_i=0}b_i^2$.
Centring is an $\ell_2$ projection for every $\tau_\mu$, so
$\sum_{i:A_i=0}b_i^2\le\sum_{i:A_i=0}\{\widehat w_{\tau_e}(Z_i)-w_0^\circ(Z_i)\}^2$
regardless of $\tau_\mu$, and, evaluated at $\tau_e=\widehat\tau_e$,
Hypothesis~\ref{h:wclass} together with Young's inequality bounds the
right side by $O_p(D_{w,n}^2+r_n^2)$ (as derived in
Hypothesis~\ref{h:wclass}'s own statement). The partition $\widehat\tau_\mu$
actually selected is, however, data-dependent: making the concentration
step below hold at the realised pair $(\widehat\tau_e,\widehat\tau_\mu)$,
rather than only at a pair fixed in advance, requires the same union
argument over the $s_n$-cardinality class used throughout (U3), which costs
an extra factor $r_n=\sqrt{s_n/n}$ relative to the fixed-pair rate. This
gives, uniformly over both indices,
\[
  |T_\varepsilon|=O_p\bigl\{(D_{w,n}+r_n)\,r_n\bigr\}.
\]

For $T_a$, Cauchy--Schwarz for the empirical control measure gives
$|T_a|\le\{P_n[(1-A)b^2]\}^{1/2}\{P_n[(1-A)a_{\mu,\tau_\mu}^2]\}^{1/2}$. The
first factor is bounded exactly as above (again
$\le n^{-1}\sum_{i:A_i=0}\{\widehat w_{\tau_e}-w_0^\circ\}^2(Z_i)
=O_p(D_{w,n}^2+r_n^2)$ at $\tau_e=\widehat\tau_e$, regardless of
$\tau_\mu$); the second factor is bounded by Hypothesis~\ref{h:aclass}'s
second display, uniformly over the deterministic family $\{a_{\mu,\tau}\}$,
at $P[(1-A)a_{\mu,\tau_\mu}^2]+O_p(\|a_{\mu,\tau_\mu}\|_\mu r_n+r_n^2)
\le D_{\mu,n}^2+O_p(D_{\mu,n}r_n+r_n^2)=O_p(D_{\mu,n}^2+r_n^2)$, using
$\|a_{\mu,\tau_\mu}\|_\mu\le D_{\mu,n}$ by $\nu$-projection optimality and
Young's inequality for the last step. Multiplying the two \emph{marginal}
bounds -- no joint union over pairs of trees is needed here, only the two
one-class Hypotheses~\ref{h:wclass}
and~\ref{h:aclass} -- gives
$|T_a|=O_p\{D_{w,n}D_{\mu,n}+(D_{w,n}+D_{\mu,n})r_n+r_n^2\}$. Thus
\[
 |T|
 =O_p\!\left[
 (D_{w,n}+r_n)\,r_n
 +D_{w,n}D_{\mu,n}
 +(D_{w,n}+D_{\mu,n})r_n+r_n^2
 \right].
\]
After multiplication by $\sqrt n$: the $T_\varepsilon$ term gives
$\sqrt n(D_{w,n}+r_n)r_n=\sqrt nD_{w,n}r_n+\sqrt nr_n^2
=O_p\{\sqrt{h_ns_n}+s_n/\sqrt n\}=o_p(1)$ by the same computation as
Corollary~\ref{cor:grid-crossfit}'s; the product term $D_{w,n}D_{\mu,n}$ is
$o_p(n^{-1/2})$ by that same corollary; the remaining cross term
$\sqrt n(D_{w,n}+D_{\mu,n})r_n=O_p\{\sqrt{h_ns_n}+s_n/\sqrt n\}=o_p(1)$
likewise; and $\sqrt nr_n^2=s_n/\sqrt n\to0$. Therefore $\sqrt nT=o_p(1)$.

\emph{Step 3: full-sample asymptotic linearity.} Lemma~\ref{lem:skeleton}
is an exact algebraic identity and does not require structural sparsity.
Hence $\widehat\pi(\widehat\theta_n-\theta_0)=\pi P_n\psi_0+U-T$. Steps 1
and 2 give $\sqrt n(U-T)=o_p(1)$. Also $\widehat\pi-\pi=O_p(n^{-1/2})$ and
$\sqrt nP_n\psi_0=O_p(1)$, so
$(\pi/\widehat\pi-1)\sqrt nP_n\psi_0=o_p(1)$, and therefore
$\sqrt n(\widehat\theta_n-\theta_0)=\sqrt nP_n\psi_0+o_p(1)$. Under the
$Z$-indexed sampling, identification, and moment conditions of
Assumption~\ref{ass:mesh} (the analogues of
Assumptions~\ref{ass:iid}, \ref{ass:causal}, and~\ref{ass:moments} at the
level of $Z$), the summands $\psi_0(O_i)$ are i.i.d., mean zero, with
variance $V\in(0,\infty)$ (the continuum quantity declared in
Assumption~\ref{ass:mesh}). The classical central limit theorem and
Slutsky's theorem give
$\sqrt n(\widehat\theta_n-\theta_0)\rightsquigarrow N(0,V)$.

\emph{Step 4: plug-in variance under vanishing misspecification.} Let
$\widehat g_n$ denote the bracketed term in \eqref{eq:vhat-def} and
$g_0=\pi\psi_0$. Bounded propensity weights and bounded outcomes imply
$(\widehat g_n-g_0)^2\le C[(\widehat e_n-\ez^\circ)^2
+(\widehat\mu_n-\mz^\circ)^2+(\widehat\theta_n-\theta_0)^2]$ for a constant
$C$ independent of $n$. The localized empirical-norm control,
Lemma~\ref{lem:grid-oracle}, and Step 3 therefore give
$P_n(\widehat g_n-g_0)^2=o_p(1)$. Cauchy--Schwarz yields
$|P_n\widehat g_n^2-P_ng_0^2|
\le\{P_n(\widehat g_n-g_0)^2\}^{1/2}\{P_n(\widehat g_n+g_0)^2\}^{1/2}
=o_p(1)$, because the first factor is $o_p(1)$ and the second is $O_p(1)$.
The weak law gives $P_ng_0^2\to_pE(g_0^2)=\pi^2V$ and
$\widehat\pi\to_p\pi>0$. Consequently
$\widehat V_n=\widehat\pi^{-2}P_n\widehat g_n^2\to_pV$. Studentization and
Slutsky give the asserted Wald coverage. Remark~\ref{rem:honest}'s
fixed-misspecification failure does not contradict this: that expansion's
first-order discrepancy is linear in $a=\mu_*-\mz$ and $b=w_*-w_0$; here
$\|a\|_2+\|b\|_2=o_p(1)$, so that discrepancy and its quadratic remainder
both vanish, rather than the first-order term dominating a vanishing
remainder.

\emph{Step 5: regularity.} Step 3 established
$\sqrt n(\widehat\theta_n-\theta_0)=\sqrt nP_n\psi_0+o_p(1)$ at the fixed
continuum law $P_0$, for $\psi_0$ Lemma~\ref{lem:tangent}'s canonical
gradient. If, in addition, $P_0$ satisfies the interiority condition of
Lemma~\ref{lem:tangent}, this is exactly the hypothesis of
Theorem~\ref{thm:regular}, which applies directly and gives regularity and
attainment of $V$.
\end{proof}

\paragraph{Post-proof audit.}
\emph{Assumptions:} unlike the companion document's main efficiency
theorem, this theorem does not use structural sparsity, the population
margin lemma, the selection-consistency lemma, or a selection event of the
form $\{\widehat\tau_j\in\cS_j\}$.
\emph{Dependencies:} Lemma~\ref{lem:skeleton}, Proposition~\ref{prop:bilinear},
Lemma~\ref{lem:clip}, Lemma~\ref{lem:normequiv}, and the oracle inequality
are reused; the off-sparsity bounds replacing the companion document's
leaf-value and cross-term lemmas are new.
\emph{Rates:} every full-sample remainder reduces to one of $\sqrt{h_ns_n}$,
$s_n/\sqrt n$, $\sqrt n h_n$, or a vanishing nuisance norm; all converge to
zero under~\eqref{eq:grid-rootn}.
\emph{Edge cases:} weak overlap destroys bounded weights; unbounded
outcomes require a replacement tail condition; loss of comparator leaf mass
destroys feasibility; and U1 without U1$^+$ does not ensure root-$n$
negligibility.
\emph{Tightness:} the entropy restriction may be stronger than necessary if
a sharper local-complexity argument is available.
\emph{Failure localization:} the least certain ingredient is
Hypothesis~\ref{h:refit} and its accompanying rate condition $s_n=o(m_n)$:
unlike Hypotheses~\ref{h:wclass} and~\ref{h:aclass} (each a marginal,
single-class localized concentration statement of the same form used
throughout (U3)), Hypothesis~\ref{h:refit}'s leaf-refit concentration
additionally requires the SELECTED tree's leaves to be occupied densely
enough, relative to $s_n$, for the per-leaf refit noise to concentrate; this
condition is new relative to Assumption~\ref{ass:grid-inference} and is not
independently derived here. A proof that supplies only an additive
$O_p(r_n)$ bound for squared norms, rather than the relative
$O_p(D^2+r_n^2)$ bounds in Hypotheses~\ref{h:wclass}--\ref{h:aclass}, is
insufficient at the root-$n$ scale. No grid, partition, projection, or
nuisance fit is treated as invariant across $n$.

% ===========================================================================
\section{Discussion and open questions}\label{sec:discussion}

\begin{enumerate}[label=(\alph*),leftmargin=*]

\item The oracle inequality of Lemma~\ref{lem:grid-oracle} delivers
$O_p(\delta^*_{j,n}+\sqrt{\log\card{\cT_{\bar L,n}}/n})$, not
$O_p(\delta^*_{j,n}+n^{-1/2})$, because (U1) implies a genuinely growing
class. The stronger entropy condition $s_n=o(\sqrt n)$ was added for
Results A and B accordingly; (U1) alone is enough for consistency but not
for the DML product condition that Result A's proof needs.

\item The companion document's condition $\bar Lm_n\lesssim n$, with
$\bar L$ fixed, does not by itself imply $m_n/n\to0$. The latter is stated
explicitly in (U2) because it is needed to make the positive-mass
comparator asymptotically feasible.

\item (U2) is derived from~\eqref{eq:mesh-slab} after deleting null
continuum leaves. It does not provide a mass floor for every candidate in
$\cT_{\bar L,n}$, nor is such a floor claimed.

\item Lemma~\ref{lem:mesh-approx}'s grid-alignment error $\delta^*_{j,n}$ is
now defined directly against the continuum truth
(\eqref{eq:mesh-delta}, revised): since every $\tau\in\cT_{\bar L,n}$'s
leaves are unions of $X_n$ cells, the tower property gives
$\Pi_\tau\ezn=\Pi_\tau\ez^\circ$ and $\Pi^\nu_\tau\mzn=\Pi^\nu_\tau\mz^\circ$
identically, so there is no residual discretisation gap between the
"discretized-truth" and "continuum-truth" readings of $\delta^*_{j,n}$ to
reconcile. Consequently Lemma~\ref{lem:grid-oracle}'s conclusion is stated
as a single equality, $O_p(\sqrt{h_n}+r_n)=O_p(\delta^*_{j,n}+r_n)$, with no
separate "sharper" case and no additional oracle leaf-mass property beyond
(U2).

\item The identification question left open in Assumption~\ref{ass:mesh}
is resolved directly in that assumption's statement: the causal model,
moment bound, and semiparametric variance bound are all stated at the level
of $Z$, and $\theta_0$ is declared there as the continuum-identified
target, $\theta_0=\pi^{-1}\E[\ez^\circ(m_1^\circ-\mz^\circ)]$. The trees are
restricted to functions of $X_n=Q_n(Z)$, but their errors are measured
against the continuum nuisances throughout Section~\ref{sec:grid-inference}.
If only $X_n$ is observed (rather than $Z$), an additional assumption
establishing ignorability given $X_n$ and equality of the resulting row-$n$
target with $\theta_0$ would be required; it does not follow from
Assumption~\ref{ass:mesh} and is not addressed here.

\item Result A (Corollary~\ref{cor:grid-crossfit}) is a genuine corollary
of Theorem~\ref{thm:crossfit}. Result A$'$ (Corollary~\ref{cor:grid-width})
uses Theorem~\ref{thm:anchor} only for validity and makes no efficiency
claim. Result B (Theorem~\ref{thm:grid-fullsample}) is independent of the
companion document's main efficiency theorem: it replaces, rather than
extends, that theorem's sparsity-dependent selection-event proof, using a
different technique throughout.

\item Theorem~\ref{thm:grid-fullsample}'s hypotheses are now stated as
three marginal, single-class conditions
(Hypotheses~\ref{h:wclass}--\ref{h:refit}) rather than one joint hypothesis
over pairs of trees: neither the $T_\varepsilon$ nor the $T_a$ bound in the
proof's Step 2 in fact requires joint control over $\tau_e\times\tau_\mu$
simultaneously, only marginal control over each index separately (with the
data-dependence of the \emph{other} index handled by an $\ell_2$-projection
identity that holds pointwise, for every candidate, not probabilistically).
The one place a genuine uniformity-over-\emph{both}-indices statement is
needed is the $T_\varepsilon$ concentration step itself (not the $T_a$
bound, which was this document's earlier, and misplaced, self-diagnosis);
Step 2 now makes that union explicit and prices it at the same $r_n$ rate
used throughout (U3).

\item Hypothesis~\ref{h:refit} (Step 1's leaf-refit term) is the one
genuinely new primitive condition added in this revision, carrying the
additional rate requirement $s_n=o(m_n)$: the per-leaf refit noise
concentrates only if the class's log-cardinality grows strictly slower than
the feasible-set's leaf-occupancy floor. This condition is stated, not
derived from anything more primitive within this document, and is the
place a subsequent reviewer should look first. It is, however, compatible
with the other restrictions on $m_n$: under (U1)'s leading case
$s_n=O(\log n)$, any $m_n$ with $\log n\ll m_n\ll n$ satisfies
$s_n=o(m_n)$, $m_n/n\to0$ (Assumption~\ref{ass:grid-inference}(U2)), and
$\bar Lm_n\lesssim n$ (Assumption~\ref{ass:construct}(d)) simultaneously --
for instance $m_n\asymp\sqrt n$.

\end{enumerate}

% ===========================================================================
\begin{thebibliography}{9}

\bibitem[Armstrong \& Kolesár(2018)]{armstrong2018} Armstrong, T. B. and
  Kolesár, M. (2018). Optimal inference in a class of regression models.
  \emph{Econometrica} 86, 655--683.

\bibitem[Barrow \& Smith(1978)]{barrowsmith1978} Barrow, D. L. and Smith,
  P. W. (1978). Asymptotic properties of best $L_2[0,1]$ approximation by
  splines with variable knots. \emph{Quarterly of Applied Mathematics} 36.
  \url{https://doi.org/10.1090/qam/508773}.

\bibitem[Chernozhukov et al.(2018)]{chernozhukov2018} Chernozhukov, V.,
  Chetverikov, D., Demirer, M., Duflo, E., Hansen, C., Newey, W. and Robins,
  J. (2018). Double/debiased machine learning for treatment and structural
  parameters. \emph{The Econometrics Journal} 21, C1--C68.

\bibitem[de Boor(1973)]{deboor1973} de Boor, C. (1973). Good approximation
  by splines with variable knots. In A. Meir and A. Sharma (eds.),
  \emph{Spline Functions and Approximation Theory}, ISNM 21, 57--72.
  Birkh\"auser, Basel. \url{https://doi.org/10.1007/978-3-0348-5979-0_3}.

\bibitem[Lugosi \& Nobel(1996)]{lugosinobel1996} Lugosi, G. and Nobel, A.
  (1996). Consistency of data-driven histogram methods for density
  estimation and classification. \emph{The Annals of Statistics} 24(2).
  \url{https://doi.org/10.1214/aos/1032894460}.

\bibitem[Nobel(1996)]{nobel1996} Nobel, A. B. (1996). Histogram regression
  estimation using data-dependent partitions. \emph{The Annals of
  Statistics} 24(3). \url{https://doi.org/10.1214/aos/1032526958}.

\end{thebibliography}

\end{document}
